% paper29-repair-semantics.tex — Sheaf-Guided Program Repair: From Obstruction to Fix
% Paper 29 of the Judgment Geometry series.
% Compile: pdflatex -interaction=nonstopmode paper29-repair-semantics.tex
\documentclass[11pt]{article}
\usepackage{lmodern}
% jugeo-common.tex — shared preamble for all Judgment Geometry papers
% Include via: % jugeo-common.tex — shared preamble for all Judgment Geometry papers
% Include via: % jugeo-common.tex — shared preamble for all Judgment Geometry papers
% Include via: \input{jugeo-common}

% ── Packages ────────────────────────────────────────────────────────────
\usepackage{amsmath,amssymb,amsthm,mathtools}
\usepackage{tikz-cd}
\usepackage{listings}
\usepackage{xcolor}
\usepackage{xspace}
\usepackage{booktabs}
\usepackage{multirow}
\usepackage{enumitem}
\usepackage[expansion=false]{microtype}
\usepackage{hyperref}
\usepackage{cleveref}
% \usepackage{thmtools}  % disabled: not available in this TeX installation
\usepackage{float}
\usepackage{caption}
\usepackage{subcaption}
\usepackage{tabularx}
\usepackage{stmaryrd}       % \llbracket, \rrbracket
\usepackage{bbm}            % \mathbbm{1}

% ── Page geometry ───────────────────────────────────────────────────────
\usepackage[margin=1in]{geometry}

% ── Theorem environments ────────────────────────────────────────────────
\theoremstyle{plain}
\newtheorem{theorem}{Theorem}[section]
\newtheorem{lemma}[theorem]{Lemma}
\newtheorem{proposition}[theorem]{Proposition}
\newtheorem{corollary}[theorem]{Corollary}

\theoremstyle{definition}
\newtheorem{definition}[theorem]{Definition}
\newtheorem{example}[theorem]{Example}
\newtheorem{construction}[theorem]{Construction}

\theoremstyle{remark}
\newtheorem{remark}[theorem]{Remark}
\newtheorem{notation}[theorem]{Notation}

% ── Python listings style ──────────────────────────────────────────────
\definecolor{jugeoBlue}{HTML}{2563EB}
\definecolor{jugeoGreen}{HTML}{059669}
\definecolor{jugeoRed}{HTML}{DC2626}
\definecolor{jugeoPurple}{HTML}{7C3AED}
\definecolor{jugeoGray}{HTML}{6B7280}
\definecolor{jugeoBg}{HTML}{F8FAFC}

\lstdefinestyle{jugeo-python}{
  language=Python,
  basicstyle=\ttfamily\small,
  keywordstyle=\color{jugeoBlue}\bfseries,
  stringstyle=\color{jugeoGreen},
  commentstyle=\color{jugeoGray}\itshape,
  numberstyle=\tiny\color{jugeoGray},
  backgroundcolor=\color{jugeoBg},
  numbers=left,
  numbersep=6pt,
  frame=single,
  framerule=0.4pt,
  rulecolor=\color{jugeoGray!40},
  breaklines=true,
  breakatwhitespace=true,
  showstringspaces=false,
  tabsize=4,
  xleftmargin=1.5em,
  framexleftmargin=1.5em,
  morekeywords={dataclass,frozen,field,Enum,IntEnum,auto,Any,
                Mapping,Sequence,Optional,match,case},
  literate={→}{{$\to$}}1 {←}{{$\leftarrow$}}1
           {≤}{{$\leq$}}1 {≥}{{$\geq$}}1 {≠}{{$\neq$}}1
           {∈}{{$\in$}}1 {∀}{{$\forall$}}1 {∃}{{$\exists$}}1
           {⊕}{{$\oplus$}}1 {⊖}{{$\ominus$}}1,
  escapeinside={(*@}{@*)},
}
\lstset{style=jugeo-python}

\lstdefinestyle{jugeo-output}{
  basicstyle=\ttfamily\small,
  backgroundcolor=\color{jugeoBg},
  frame=single,
  framerule=0.4pt,
  rulecolor=\color{jugeoGray!40},
  breaklines=true,
  showstringspaces=false,
  xleftmargin=1.5em,
  framexleftmargin=0.5em,
  numbers=none,
}

% ── JuGeo-specific macros ──────────────────────────────────────────────

% Judgment 8-tuple
\newcommand{\J}{\mathcal{J}}
\newcommand{\Jud}[8]{(#1,\,#2,\,#3,\,#4,\,#5,\,#6,\,#7,\,#8)}
\newcommand{\JudShort}[1]{\J_{#1}}

% Components
\newcommand{\coord}{\mathsf{c}}            % coordinate
\newcommand{\prop}{\varphi}                 % proposition
\newcommand{\carrier}{\mathcal{A}}          % carrier type
\newcommand{\evid}{\mathcal{E}}             % evidence bundle
\newcommand{\oblig}{\mathcal{O}}            % obligations
\newcommand{\obstr}{\mathcal{B}}            % obstructions
\newcommand{\trust}{\mathcal{T}}            % trust annotation
\newcommand{\prov}{\Pi}                     % provenance

% Trust algebra
\newcommand{\Talg}{\mathfrak{T}}            % trust ordered algebra
\newcommand{\Eadm}{\mathcal{E}_{\mathrm{adm}}}
\newcommand{\tleq}{\preceq}                 % trust ordering
\newcommand{\tjoin}{\oplus}                 % conservative join
\newcommand{\tatten}{\ominus}               % attenuation
\newcommand{\tpromote}{\uparrow_\pi}        % promotion
\newcommand{\tdemote}{\downarrow_\chi}       % demotion

% Trust tiers
\newcommand{\Tcontra}{\ensuremath{\mathsf{CONTRADICTED}}}
\newcommand{\Tunver}{\ensuremath{\mathsf{UNVERIFIED}}}
\newcommand{\Tcopilot}{\ensuremath{\mathsf{COPILOT}}}
\newcommand{\Toracle}{\ensuremath{\mathsf{ORACLE}}}
\newcommand{\Truntime}{\ensuremath{\mathsf{RUNTIME}}}
\newcommand{\Tsolver}{\ensuremath{\mathsf{SOLVER}}}
\newcommand{\Tproof}{\ensuremath{\mathsf{PROOF}}}

% Site and geometry
\newcommand{\Site}{\mathcal{S}}             % semantic site
\newcommand{\Cov}{\mathrm{Cov}}            % covering family
\newcommand{\Ob}{\mathrm{Ob}}              % objects
\newcommand{\Mor}{\mathrm{Mor}}            % morphisms
\newcommand{\res}{\mathrm{res}}            % restriction
\newcommand{\Cech}{\check{C}}              % Čech complex
\newcommand{\Hcoh}[1]{H^{#1}}             % cohomology group

% Descent
\newcommand{\descent}{\mathrm{desc}}
\newcommand{\glue}{\mathrm{glue}}
\newcommand{\overlap}[2]{#1 \cap #2}

% Semantic moves
\newcommand{\VERIFY}{\textsc{Verify}}
\newcommand{\CONSTRUCT}{\textsc{Construct}}
\newcommand{\NORMALIZE}{\textsc{Normalize}}
\newcommand{\GLUE}{\textsc{Glue}}
\newcommand{\RESTRICT}{\textsc{Restrict}}
\newcommand{\TRANSPORT}{\textsc{Transport}}
\newcommand{\REPAIR}{\textsc{Repair}}
\newcommand{\PROMOTE}{\textsc{Promote}}
\newcommand{\CHALLENGE}{\textsc{Challenge}}
\newcommand{\ESCALATE}{\textsc{Escalate}}

% Fragments
\newcommand{\QFLIA}{\texttt{QF\_LIA}}
\newcommand{\QFLRA}{\texttt{QF\_LRA}}
\newcommand{\QFBV}{\texttt{QF\_BV}}
\newcommand{\QFUF}{\texttt{QF\_UF}}

% Misc
\newcommand{\Python}{\textsc{Python}}
\newcommand{\JuGeo}{\textsc{JuGeo}}
\newcommand{\LEAN}{\textsc{Lean}\,4}
\newcommand{\Fstar}{F$^{\star}$}
\newcommand{\Coq}{\textsc{Coq}}
\newcommand{\Dafny}{\textsc{Dafny}}

% Common shortcuts
\newcommand{\ie}{i.e.\@\xspace}
\newcommand{\eg}{e.g.\@\xspace}
\newcommand{\cf}{cf.\@\xspace}
\newcommand{\wrt}{w.r.t.\@\xspace}

% Evaluation shortcuts
\newcommand{\numcases}{300}
\newcommand{\numspec}{100}
\newcommand{\numeq}{100}
\newcommand{\numbug}{100}
\newcommand{\medtime}{6.5\,ms}
\newcommand{\totaltime}{1.949\,s}


% ── Packages ────────────────────────────────────────────────────────────
\usepackage{amsmath,amssymb,amsthm,mathtools}
\usepackage{tikz-cd}
\usepackage{listings}
\usepackage{xcolor}
\usepackage{xspace}
\usepackage{booktabs}
\usepackage{multirow}
\usepackage{enumitem}
\usepackage[expansion=false]{microtype}
\usepackage{hyperref}
\usepackage{cleveref}
% \usepackage{thmtools}  % disabled: not available in this TeX installation
\usepackage{float}
\usepackage{caption}
\usepackage{subcaption}
\usepackage{tabularx}
\usepackage{stmaryrd}       % \llbracket, \rrbracket
\usepackage{bbm}            % \mathbbm{1}

% ── Page geometry ───────────────────────────────────────────────────────
\usepackage[margin=1in]{geometry}

% ── Theorem environments ────────────────────────────────────────────────
\theoremstyle{plain}
\newtheorem{theorem}{Theorem}[section]
\newtheorem{lemma}[theorem]{Lemma}
\newtheorem{proposition}[theorem]{Proposition}
\newtheorem{corollary}[theorem]{Corollary}

\theoremstyle{definition}
\newtheorem{definition}[theorem]{Definition}
\newtheorem{example}[theorem]{Example}
\newtheorem{construction}[theorem]{Construction}

\theoremstyle{remark}
\newtheorem{remark}[theorem]{Remark}
\newtheorem{notation}[theorem]{Notation}

% ── Python listings style ──────────────────────────────────────────────
\definecolor{jugeoBlue}{HTML}{2563EB}
\definecolor{jugeoGreen}{HTML}{059669}
\definecolor{jugeoRed}{HTML}{DC2626}
\definecolor{jugeoPurple}{HTML}{7C3AED}
\definecolor{jugeoGray}{HTML}{6B7280}
\definecolor{jugeoBg}{HTML}{F8FAFC}

\lstdefinestyle{jugeo-python}{
  language=Python,
  basicstyle=\ttfamily\small,
  keywordstyle=\color{jugeoBlue}\bfseries,
  stringstyle=\color{jugeoGreen},
  commentstyle=\color{jugeoGray}\itshape,
  numberstyle=\tiny\color{jugeoGray},
  backgroundcolor=\color{jugeoBg},
  numbers=left,
  numbersep=6pt,
  frame=single,
  framerule=0.4pt,
  rulecolor=\color{jugeoGray!40},
  breaklines=true,
  breakatwhitespace=true,
  showstringspaces=false,
  tabsize=4,
  xleftmargin=1.5em,
  framexleftmargin=1.5em,
  morekeywords={dataclass,frozen,field,Enum,IntEnum,auto,Any,
                Mapping,Sequence,Optional,match,case},
  literate={→}{{$\to$}}1 {←}{{$\leftarrow$}}1
           {≤}{{$\leq$}}1 {≥}{{$\geq$}}1 {≠}{{$\neq$}}1
           {∈}{{$\in$}}1 {∀}{{$\forall$}}1 {∃}{{$\exists$}}1
           {⊕}{{$\oplus$}}1 {⊖}{{$\ominus$}}1,
  escapeinside={(*@}{@*)},
}
\lstset{style=jugeo-python}

\lstdefinestyle{jugeo-output}{
  basicstyle=\ttfamily\small,
  backgroundcolor=\color{jugeoBg},
  frame=single,
  framerule=0.4pt,
  rulecolor=\color{jugeoGray!40},
  breaklines=true,
  showstringspaces=false,
  xleftmargin=1.5em,
  framexleftmargin=0.5em,
  numbers=none,
}

% ── JuGeo-specific macros ──────────────────────────────────────────────

% Judgment 8-tuple
\newcommand{\J}{\mathcal{J}}
\newcommand{\Jud}[8]{(#1,\,#2,\,#3,\,#4,\,#5,\,#6,\,#7,\,#8)}
\newcommand{\JudShort}[1]{\J_{#1}}

% Components
\newcommand{\coord}{\mathsf{c}}            % coordinate
\newcommand{\prop}{\varphi}                 % proposition
\newcommand{\carrier}{\mathcal{A}}          % carrier type
\newcommand{\evid}{\mathcal{E}}             % evidence bundle
\newcommand{\oblig}{\mathcal{O}}            % obligations
\newcommand{\obstr}{\mathcal{B}}            % obstructions
\newcommand{\trust}{\mathcal{T}}            % trust annotation
\newcommand{\prov}{\Pi}                     % provenance

% Trust algebra
\newcommand{\Talg}{\mathfrak{T}}            % trust ordered algebra
\newcommand{\Eadm}{\mathcal{E}_{\mathrm{adm}}}
\newcommand{\tleq}{\preceq}                 % trust ordering
\newcommand{\tjoin}{\oplus}                 % conservative join
\newcommand{\tatten}{\ominus}               % attenuation
\newcommand{\tpromote}{\uparrow_\pi}        % promotion
\newcommand{\tdemote}{\downarrow_\chi}       % demotion

% Trust tiers
\newcommand{\Tcontra}{\ensuremath{\mathsf{CONTRADICTED}}}
\newcommand{\Tunver}{\ensuremath{\mathsf{UNVERIFIED}}}
\newcommand{\Tcopilot}{\ensuremath{\mathsf{COPILOT}}}
\newcommand{\Toracle}{\ensuremath{\mathsf{ORACLE}}}
\newcommand{\Truntime}{\ensuremath{\mathsf{RUNTIME}}}
\newcommand{\Tsolver}{\ensuremath{\mathsf{SOLVER}}}
\newcommand{\Tproof}{\ensuremath{\mathsf{PROOF}}}

% Site and geometry
\newcommand{\Site}{\mathcal{S}}             % semantic site
\newcommand{\Cov}{\mathrm{Cov}}            % covering family
\newcommand{\Ob}{\mathrm{Ob}}              % objects
\newcommand{\Mor}{\mathrm{Mor}}            % morphisms
\newcommand{\res}{\mathrm{res}}            % restriction
\newcommand{\Cech}{\check{C}}              % Čech complex
\newcommand{\Hcoh}[1]{H^{#1}}             % cohomology group

% Descent
\newcommand{\descent}{\mathrm{desc}}
\newcommand{\glue}{\mathrm{glue}}
\newcommand{\overlap}[2]{#1 \cap #2}

% Semantic moves
\newcommand{\VERIFY}{\textsc{Verify}}
\newcommand{\CONSTRUCT}{\textsc{Construct}}
\newcommand{\NORMALIZE}{\textsc{Normalize}}
\newcommand{\GLUE}{\textsc{Glue}}
\newcommand{\RESTRICT}{\textsc{Restrict}}
\newcommand{\TRANSPORT}{\textsc{Transport}}
\newcommand{\REPAIR}{\textsc{Repair}}
\newcommand{\PROMOTE}{\textsc{Promote}}
\newcommand{\CHALLENGE}{\textsc{Challenge}}
\newcommand{\ESCALATE}{\textsc{Escalate}}

% Fragments
\newcommand{\QFLIA}{\texttt{QF\_LIA}}
\newcommand{\QFLRA}{\texttt{QF\_LRA}}
\newcommand{\QFBV}{\texttt{QF\_BV}}
\newcommand{\QFUF}{\texttt{QF\_UF}}

% Misc
\newcommand{\Python}{\textsc{Python}}
\newcommand{\JuGeo}{\textsc{JuGeo}}
\newcommand{\LEAN}{\textsc{Lean}\,4}
\newcommand{\Fstar}{F$^{\star}$}
\newcommand{\Coq}{\textsc{Coq}}
\newcommand{\Dafny}{\textsc{Dafny}}

% Common shortcuts
\newcommand{\ie}{i.e.\@\xspace}
\newcommand{\eg}{e.g.\@\xspace}
\newcommand{\cf}{cf.\@\xspace}
\newcommand{\wrt}{w.r.t.\@\xspace}

% Evaluation shortcuts
\newcommand{\numcases}{300}
\newcommand{\numspec}{100}
\newcommand{\numeq}{100}
\newcommand{\numbug}{100}
\newcommand{\medtime}{6.5\,ms}
\newcommand{\totaltime}{1.949\,s}


% ── Packages ────────────────────────────────────────────────────────────
\usepackage{amsmath,amssymb,amsthm,mathtools}
\usepackage{tikz-cd}
\usepackage{listings}
\usepackage{xcolor}
\usepackage{xspace}
\usepackage{booktabs}
\usepackage{multirow}
\usepackage{enumitem}
\usepackage[expansion=false]{microtype}
\usepackage{hyperref}
\usepackage{cleveref}
% \usepackage{thmtools}  % disabled: not available in this TeX installation
\usepackage{float}
\usepackage{caption}
\usepackage{subcaption}
\usepackage{tabularx}
\usepackage{stmaryrd}       % \llbracket, \rrbracket
\usepackage{bbm}            % \mathbbm{1}

% ── Page geometry ───────────────────────────────────────────────────────
\usepackage[margin=1in]{geometry}

% ── Theorem environments ────────────────────────────────────────────────
\theoremstyle{plain}
\newtheorem{theorem}{Theorem}[section]
\newtheorem{lemma}[theorem]{Lemma}
\newtheorem{proposition}[theorem]{Proposition}
\newtheorem{corollary}[theorem]{Corollary}

\theoremstyle{definition}
\newtheorem{definition}[theorem]{Definition}
\newtheorem{example}[theorem]{Example}
\newtheorem{construction}[theorem]{Construction}

\theoremstyle{remark}
\newtheorem{remark}[theorem]{Remark}
\newtheorem{notation}[theorem]{Notation}

% ── Python listings style ──────────────────────────────────────────────
\definecolor{jugeoBlue}{HTML}{2563EB}
\definecolor{jugeoGreen}{HTML}{059669}
\definecolor{jugeoRed}{HTML}{DC2626}
\definecolor{jugeoPurple}{HTML}{7C3AED}
\definecolor{jugeoGray}{HTML}{6B7280}
\definecolor{jugeoBg}{HTML}{F8FAFC}

\lstdefinestyle{jugeo-python}{
  language=Python,
  basicstyle=\ttfamily\small,
  keywordstyle=\color{jugeoBlue}\bfseries,
  stringstyle=\color{jugeoGreen},
  commentstyle=\color{jugeoGray}\itshape,
  numberstyle=\tiny\color{jugeoGray},
  backgroundcolor=\color{jugeoBg},
  numbers=left,
  numbersep=6pt,
  frame=single,
  framerule=0.4pt,
  rulecolor=\color{jugeoGray!40},
  breaklines=true,
  breakatwhitespace=true,
  showstringspaces=false,
  tabsize=4,
  xleftmargin=1.5em,
  framexleftmargin=1.5em,
  morekeywords={dataclass,frozen,field,Enum,IntEnum,auto,Any,
                Mapping,Sequence,Optional,match,case},
  literate={→}{{$\to$}}1 {←}{{$\leftarrow$}}1
           {≤}{{$\leq$}}1 {≥}{{$\geq$}}1 {≠}{{$\neq$}}1
           {∈}{{$\in$}}1 {∀}{{$\forall$}}1 {∃}{{$\exists$}}1
           {⊕}{{$\oplus$}}1 {⊖}{{$\ominus$}}1,
  escapeinside={(*@}{@*)},
}
\lstset{style=jugeo-python}

\lstdefinestyle{jugeo-output}{
  basicstyle=\ttfamily\small,
  backgroundcolor=\color{jugeoBg},
  frame=single,
  framerule=0.4pt,
  rulecolor=\color{jugeoGray!40},
  breaklines=true,
  showstringspaces=false,
  xleftmargin=1.5em,
  framexleftmargin=0.5em,
  numbers=none,
}

% ── JuGeo-specific macros ──────────────────────────────────────────────

% Judgment 8-tuple
\newcommand{\J}{\mathcal{J}}
\newcommand{\Jud}[8]{(#1,\,#2,\,#3,\,#4,\,#5,\,#6,\,#7,\,#8)}
\newcommand{\JudShort}[1]{\J_{#1}}

% Components
\newcommand{\coord}{\mathsf{c}}            % coordinate
\newcommand{\prop}{\varphi}                 % proposition
\newcommand{\carrier}{\mathcal{A}}          % carrier type
\newcommand{\evid}{\mathcal{E}}             % evidence bundle
\newcommand{\oblig}{\mathcal{O}}            % obligations
\newcommand{\obstr}{\mathcal{B}}            % obstructions
\newcommand{\trust}{\mathcal{T}}            % trust annotation
\newcommand{\prov}{\Pi}                     % provenance

% Trust algebra
\newcommand{\Talg}{\mathfrak{T}}            % trust ordered algebra
\newcommand{\Eadm}{\mathcal{E}_{\mathrm{adm}}}
\newcommand{\tleq}{\preceq}                 % trust ordering
\newcommand{\tjoin}{\oplus}                 % conservative join
\newcommand{\tatten}{\ominus}               % attenuation
\newcommand{\tpromote}{\uparrow_\pi}        % promotion
\newcommand{\tdemote}{\downarrow_\chi}       % demotion

% Trust tiers
\newcommand{\Tcontra}{\ensuremath{\mathsf{CONTRADICTED}}}
\newcommand{\Tunver}{\ensuremath{\mathsf{UNVERIFIED}}}
\newcommand{\Tcopilot}{\ensuremath{\mathsf{COPILOT}}}
\newcommand{\Toracle}{\ensuremath{\mathsf{ORACLE}}}
\newcommand{\Truntime}{\ensuremath{\mathsf{RUNTIME}}}
\newcommand{\Tsolver}{\ensuremath{\mathsf{SOLVER}}}
\newcommand{\Tproof}{\ensuremath{\mathsf{PROOF}}}

% Site and geometry
\newcommand{\Site}{\mathcal{S}}             % semantic site
\newcommand{\Cov}{\mathrm{Cov}}            % covering family
\newcommand{\Ob}{\mathrm{Ob}}              % objects
\newcommand{\Mor}{\mathrm{Mor}}            % morphisms
\newcommand{\res}{\mathrm{res}}            % restriction
\newcommand{\Cech}{\check{C}}              % Čech complex
\newcommand{\Hcoh}[1]{H^{#1}}             % cohomology group

% Descent
\newcommand{\descent}{\mathrm{desc}}
\newcommand{\glue}{\mathrm{glue}}
\newcommand{\overlap}[2]{#1 \cap #2}

% Semantic moves
\newcommand{\VERIFY}{\textsc{Verify}}
\newcommand{\CONSTRUCT}{\textsc{Construct}}
\newcommand{\NORMALIZE}{\textsc{Normalize}}
\newcommand{\GLUE}{\textsc{Glue}}
\newcommand{\RESTRICT}{\textsc{Restrict}}
\newcommand{\TRANSPORT}{\textsc{Transport}}
\newcommand{\REPAIR}{\textsc{Repair}}
\newcommand{\PROMOTE}{\textsc{Promote}}
\newcommand{\CHALLENGE}{\textsc{Challenge}}
\newcommand{\ESCALATE}{\textsc{Escalate}}

% Fragments
\newcommand{\QFLIA}{\texttt{QF\_LIA}}
\newcommand{\QFLRA}{\texttt{QF\_LRA}}
\newcommand{\QFBV}{\texttt{QF\_BV}}
\newcommand{\QFUF}{\texttt{QF\_UF}}

% Misc
\newcommand{\Python}{\textsc{Python}}
\newcommand{\JuGeo}{\textsc{JuGeo}}
\newcommand{\LEAN}{\textsc{Lean}\,4}
\newcommand{\Fstar}{F$^{\star}$}
\newcommand{\Coq}{\textsc{Coq}}
\newcommand{\Dafny}{\textsc{Dafny}}

% Common shortcuts
\newcommand{\ie}{i.e.\@\xspace}
\newcommand{\eg}{e.g.\@\xspace}
\newcommand{\cf}{cf.\@\xspace}
\newcommand{\wrt}{w.r.t.\@\xspace}

% Evaluation shortcuts
\newcommand{\numcases}{300}
\newcommand{\numspec}{100}
\newcommand{\numeq}{100}
\newcommand{\numbug}{100}
\newcommand{\medtime}{6.5\,ms}
\newcommand{\totaltime}{1.949\,s}

% experiment-data.tex — AUTO-GENERATED from real jugeo CLI runs
% DO NOT EDIT — regenerate with: python3 experiments/run_universal_metrics.py
% Generated from 30 programs across 6 domains

\newcommand{\expTotalPrograms}{30}
\newcommand{\expVerified}{30}
\newcommand{\expAccuracy}{100.0\%}
\newcommand{\expCoordsMin}{2}
\newcommand{\expCoordsMax}{10}
\newcommand{\expCoordsMean}{5.2}
\newcommand{\expCoordsMedian}{5.5}
\newcommand{\expPropsMin}{4}
\newcommand{\expPropsMax}{20}
\newcommand{\expPropsMean}{10.4}
\newcommand{\expPropsSum}{311}
\newcommand{\expPropsOkSum}{310}
\newcommand{\expTimeMin}{3.506\,s}
\newcommand{\expTimeMax}{6.714\,s}
\newcommand{\expTimeMean}{4.64\,s}
\newcommand{\expTimeMedian}{4.508\,s}
\newcommand{\expTimeTotal}{139.204\,s}
\newcommand{\expObstructionTotal}{0}
\newcommand{\expHOneZero}{30}
\newcommand{\expDomainCount}{6}
\newcommand{\expTrustCOPILOTSUGGESTED}{30}
\newcommand{\expDomainDsCount}{5}
\newcommand{\expDomainDsVerified}{5}
\newcommand{\expDomainDsAvgCoords}{7.6}
\newcommand{\expDomainDsAvgProps}{15.2}
\newcommand{\expDomainMathCount}{5}
\newcommand{\expDomainMathVerified}{5}
\newcommand{\expDomainMathAvgCoords}{4.8}
\newcommand{\expDomainMathAvgProps}{9.4}
\newcommand{\expDomainSmCount}{5}
\newcommand{\expDomainSmVerified}{5}
\newcommand{\expDomainSmAvgCoords}{7.6}
\newcommand{\expDomainSmAvgProps}{15.2}
\newcommand{\expDomainSortCount}{5}
\newcommand{\expDomainSortVerified}{5}
\newcommand{\expDomainSortAvgCoords}{2.4}
\newcommand{\expDomainSortAvgProps}{4.8}
\newcommand{\expDomainStrCount}{5}
\newcommand{\expDomainStrVerified}{5}
\newcommand{\expDomainStrAvgCoords}{3.2}
\newcommand{\expDomainStrAvgProps}{6.4}
\newcommand{\expDomainWebCount}{5}
\newcommand{\expDomainWebVerified}{5}
\newcommand{\expDomainWebAvgCoords}{5.6}
\newcommand{\expDomainWebAvgProps}{11.2}

% subsystem-data.tex — AUTO-GENERATED from real jugeo subsystem experiments
% DO NOT EDIT — regenerate with: python3 experiments/run_subsystem_experiments.py

\newcommand{\subCliTotal}{10}
\newcommand{\subCliVerified}{9}
\newcommand{\subCliAccuracy}{90.0\%}
\newcommand{\subCliMeanTime}{4.132\,s}
\newcommand{\subCliMinTime}{3.472\,s}
\newcommand{\subCliMaxTime}{5.372\,s}
\newcommand{\subCliTotalCoords}{54}
\newcommand{\subCliTotalProps}{107}
\newcommand{\subCliTotalPropsOk}{106}
\newcommand{\subSiteCalls}{90}
\newcommand{\subSiteSuccessful}{69}
\newcommand{\subSiteSuccessRate}{76.7\%}
\newcommand{\subSiteMeanTime}{0.0001\,s}
\newcommand{\subSiteTotalTime}{0.005\,s}
\newcommand{\subSpecificationsatisfactionSmallTime}{0.0003\,s}
\newcommand{\subSpecificationsatisfactionMediumTime}{0.0001\,s}
\newcommand{\subSpecificationsatisfactionLargeTime}{0.0001\,s}
\newcommand{\subInhabitantfleetSmallTime}{0.0\,s}
\newcommand{\subInhabitantfleetMediumTime}{0.0\,s}
\newcommand{\subInhabitantfleetLargeTime}{0.0\,s}
\newcommand{\subTheoremecologySmallTime}{0.0\,s}
\newcommand{\subTheoremecologyMediumTime}{0.0\,s}
\newcommand{\subTheoremecologyLargeTime}{0.0\,s}
\newcommand{\subStatespaceexplorationSmallTime}{0.0\,s}
\newcommand{\subStatespaceexplorationMediumTime}{0.0\,s}
\newcommand{\subStatespaceexplorationLargeTime}{0.0\,s}
\newcommand{\subRepairsemanticsSmallTime}{0.0\,s}
\newcommand{\subRepairsemanticsMediumTime}{0.0\,s}
\newcommand{\subRepairsemanticsLargeTime}{0.0\,s}
\newcommand{\subReplaygluingSmallTime}{0.0\,s}
\newcommand{\subReplaygluingMediumTime}{0.0\,s}
\newcommand{\subReplaygluingLargeTime}{0.0\,s}
\newcommand{\subSemanticfuturesSmallTime}{0.0\,s}
\newcommand{\subSemanticfuturesMediumTime}{0.0\,s}
\newcommand{\subSemanticfuturesLargeTime}{0.0\,s}
\newcommand{\subHypercovertreatySmallTime}{0.0\,s}
\newcommand{\subHypercovertreatyMediumTime}{0.0\,s}
\newcommand{\subHypercovertreatyLargeTime}{0.0\,s}
\newcommand{\subEncodeforsolverSmallTime}{0.0026\,s}
\newcommand{\subEncodeforsolverMediumTime}{0.0002\,s}
\newcommand{\subEncodeforsolverLargeTime}{0.0001\,s}
\newcommand{\subInterfaceroutingSmallTime}{0.0\,s}
\newcommand{\subInterfaceroutingMediumTime}{0.0\,s}
\newcommand{\subInterfaceroutingLargeTime}{0.0\,s}
\newcommand{\subOrchestrateverificationSmallTime}{0.0002\,s}
\newcommand{\subOrchestrateverificationMediumTime}{0.0\,s}
\newcommand{\subOrchestrateverificationLargeTime}{0.0\,s}
\newcommand{\subRunfulldescentSmallTime}{0.0\,s}
\newcommand{\subRunfulldescentMediumTime}{0.0\,s}
\newcommand{\subRunfulldescentLargeTime}{0.0\,s}
\newcommand{\subKernellifecycleSmallTime}{0.0\,s}
\newcommand{\subKernellifecycleMediumTime}{0.0\,s}
\newcommand{\subKernellifecycleLargeTime}{0.0\,s}
\newcommand{\subDiscoverypipelineSmallTime}{0.0\,s}
\newcommand{\subDiscoverypipelineMediumTime}{0.0\,s}
\newcommand{\subDiscoverypipelineLargeTime}{0.0\,s}
\newcommand{\subAnalogytransportSmallTime}{0.0\,s}
\newcommand{\subAnalogytransportMediumTime}{0.0\,s}
\newcommand{\subAnalogytransportLargeTime}{0.0\,s}
\newcommand{\subFormalcoresiteSmallTime}{0.0001\,s}
\newcommand{\subFormalcoresiteMediumTime}{0.0\,s}
\newcommand{\subFormalcoresiteLargeTime}{0.0\,s}
\newcommand{\subProblematlasSmallTime}{0.0\,s}
\newcommand{\subProblematlasMediumTime}{0.0\,s}
\newcommand{\subProblematlasLargeTime}{0.0\,s}
\newcommand{\subPublicalignmentSmallTime}{0.0\,s}
\newcommand{\subPublicalignmentMediumTime}{0.0\,s}
\newcommand{\subPublicalignmentLargeTime}{0.0\,s}
\newcommand{\subRegimebootstrappingSmallTime}{0.0\,s}
\newcommand{\subRegimebootstrappingMediumTime}{0.0\,s}
\newcommand{\subRegimebootstrappingLargeTime}{0.0\,s}
\newcommand{\subRelationalrefinementSmallTime}{0.0\,s}
\newcommand{\subRelationalrefinementMediumTime}{0.0\,s}
\newcommand{\subRelationalrefinementLargeTime}{0.0\,s}
\newcommand{\subSemanticclosureSmallTime}{0.0\,s}
\newcommand{\subSemanticclosureMediumTime}{0.0\,s}
\newcommand{\subSemanticclosureLargeTime}{0.0\,s}
\newcommand{\subTheoremeconomicsSmallTime}{0.0001\,s}
\newcommand{\subTheoremeconomicsMediumTime}{0.0\,s}
\newcommand{\subTheoremeconomicsLargeTime}{0.0\,s}
\newcommand{\subBenchmarksuiteSmallTime}{0.0\,s}
\newcommand{\subBenchmarksuiteMediumTime}{0.0\,s}
\newcommand{\subBenchmarksuiteLargeTime}{0.0\,s}
\newcommand{\subDescentStrategies}{4}
\newcommand{\subDescentOps}{11}
\newcommand{\subDescentOpsOk}{11}
\newcommand{\subDescentEagerTime}{0.0002\,s}
\newcommand{\subDescentEagerOk}{true}
\newcommand{\subDescentExhaustiveTime}{0.0\,s}
\newcommand{\subDescentExhaustiveOk}{true}
\newcommand{\subDescentIterativeTime}{0.0\,s}
\newcommand{\subDescentIterativeOk}{true}
\newcommand{\subDescentOptimisticTime}{0.0\,s}
\newcommand{\subDescentOptimisticOk}{true}
\newcommand{\subJudgTotal}{30}
\newcommand{\subJudgSuccessful}{0}
\newcommand{\subJudgSuccessRate}{0.0\%}
\newcommand{\subJudgMeanTimeUs}{7.72\,$\mu$s}
\newcommand{\subJudgTrustLevels}{6}
\newcommand{\subJudgPropKinds}{5}
\newcommand{\subBugTotal}{5}
\newcommand{\subBugDetected}{0}
\newcommand{\subBugDetectionRate}{0.0\%}
\newcommand{\subBugFalsePositiveRate}{0.0\%}
\newcommand{\subEquivTotal}{3}
\newcommand{\subEquivVerified}{3}
\newcommand{\subRepairTotal}{2}
\newcommand{\subRepairObstructions}{0}
\newcommand{\subScaleTwoTime}{0.0001\,s}
\newcommand{\subScaleFourTime}{0.0\,s}
\newcommand{\subScaleEightTime}{0.0\,s}
\newcommand{\subScaleTwelveTime}{0.0\,s}
\newcommand{\subScaleSixteenTime}{0.0\,s}
\newcommand{\subScaleTwentyTime}{0.0\,s}
\newcommand{\subTotalExperimentTime}{95.6\,s}

% data-paper29.tex — AUTO-GENERATED by exp29_repair_semantics.py
% DO NOT EDIT — regenerate with: python3 experiments/exp29_repair_semantics.py

\newcommand{\ppTwentynineSpecCount}{3}
\newcommand{\ppTwentynineSpecSuccessRate}{0.0\%}
\newcommand{\ppTwentynineSpecAvgRepairs}{0.0}
\newcommand{\ppTwentynineImplCount}{4}
\newcommand{\ppTwentynineImplSuccessRate}{0.0\%}
\newcommand{\ppTwentynineImplAvgRepairs}{0.0}
\newcommand{\ppTwentynineCoherenceCount}{3}
\newcommand{\ppTwentynineCoherenceSuccessRate}{0.0\%}
\newcommand{\ppTwentynineCoherenceAvgRepairs}{0.0}
\newcommand{\ppTwentynineRepairTotal}{10}
\newcommand{\ppTwentynineRepairSuccessRate}{0.0\%}
\newcommand{\ppTwentynineRepairSuccessCount}{0}
\newcommand{\ppTwentynineSpecAvgIter}{1.0}
\newcommand{\ppTwentynineImplAvgIter}{1.0}
\newcommand{\ppTwentynineCoherenceAvgIter}{1.0}
\newcommand{\ppTwentynineOverallAvgIter}{1.0}
\newcommand{\ppTwentynineFixImplCount}{0}
\newcommand{\ppTwentynineFixImplFrac}{0.0\%}
\newcommand{\ppTwentynineStrengthenPreCount}{0}
\newcommand{\ppTwentynineStrengthenPreFrac}{0.0\%}
\newcommand{\ppTwentynineRefineSpecCount}{0}
\newcommand{\ppTwentynineRefineSpecFrac}{0.0\%}
\newcommand{\ppTwentynineAddInvariantCount}{0}
\newcommand{\ppTwentynineAddInvariantFrac}{0.0\%}
\newcommand{\ppTwentynineSplitCoverCount}{0}
\newcommand{\ppTwentynineSplitCoverFrac}{0.0\%}
\newcommand{\ppTwentynineWeakenPostCount}{0}
\newcommand{\ppTwentynineWeakenPostFrac}{0.0\%}
\newcommand{\ppTwentynineTotalRepairActions}{0}
\newcommand{\ppTwentynineTotalPrograms}{10}


% ── Additional packages ────────────────────────────────────────────────
\usepackage{array}

% ── Paper-specific macros ──────────────────────────────────────────────
\newcommand{\RF}{\mathcal{F}}               % repair frontier
\newcommand{\RS}{\mathsf{RS}}               % repair strategy
\newcommand{\RO}{\mathsf{RO}}               % repair orchestrator
\newcommand{\Obs}{\mathcal{O}}              % obstruction set (overloads \oblig locally)
\newcommand{\nobst}{n_{\mathrm{obs}}}       % obstruction count
\newcommand{\apply}{\mathsf{apply}}         % apply repair step
\newcommand{\verify}{\mathsf{verify}}       % verify after repair
\newcommand{\hone}{H^1}                     % H^1 shorthand
\newcommand{\hzero}{H^0}                    % H^0 shorthand
\newcommand{\Patch}{\textsc{Patch}}
\newcommand{\Refactor}{\textsc{Refactor}}
\newcommand{\StrengthenPre}{\textsc{StrengthenPre}}
\newcommand{\WeakenPost}{\textsc{WeakenPost}}
\newcommand{\AddInvariant}{\textsc{AddInvariant}}
\newcommand{\SplitCover}{\textsc{SplitCover}}
\newcommand{\cls}[1]{[\![#1]\!]}            % cohomology class bracket

% ── Title information ─────────────────────────────────────────────────
\title{\textbf{Sheaf-Guided Program Repair:\\
       From Obstruction to Fix}\\[4pt]
       {\large Judgment Geometry Series, Paper 29}}
\author{The \JuGeo{} Research Group}
\date{2025}

\begin{document}
\maketitle

\begin{abstract}
When the sheaf descent procedure for a \JuGeo{} judgment fails, the
failure provides structured obstruction data that can in principle guide
repair.  This paper develops that idea as \emph{sheaf-guided program
repair}: a repair-strategy lattice, an apply-verify-iterate loop, and a
repair frontier that records remaining obstructions.  The current
benchmark evidence is negative rather than triumphant: across 10 repair
scenarios, the repository records a 0.0\% repair success rate, no
committed repair actions, and an average of one iteration per scenario.
The audited claim is therefore that the framework currently exposes
repair state and failure structure, not that it already repairs the
benchmark automatically.
\end{abstract}

\newpage

% ════════════════════════════════════════════════════════════════════════
\section{Introduction}
\label{sec:intro}
% ════════════════════════════════════════════════════════════════════════

Program repair has traditionally been driven by syntactic heuristics: mutate
an expression, run the test suite, keep the mutation if it passes.  Such
approaches are powerful but blind---they do not know \emph{why} a program is
wrong, only whether a particular change makes it pass a given oracle.  When
the bug has a global, semantic character (\eg a contract inconsistency that
manifests in multiple modules simultaneously) syntax-level mutation can waste
enormous effort on fixes that treat symptoms rather than causes.

\JuGeo{} offers a different starting point.  A program is represented as a
collection of \emph{local sections} over a semantic site $\Site$; correctness
is expressed as the existence of a global section that glues these local
pieces consistently (Paper~3, \cite{paper03}).  When global consistency
fails, the Čech descent procedure (Paper~3, Paper~11) produces a precise
cohomological witness: an element $[\eta] \in \hone(\mathcal{U},\mathcal{D})$
that pinpoints exactly which overlap conditions were violated, at which
coordinates, and with what algebraic structure.

This paper shows how to exploit that witness.  The central insight is:
\begin{quote}
\emph{An obstruction class is not just evidence of failure---it is a
repair roadmap.}
\end{quote}
The coordinate of each obstruction names the program region that must be
modified; the failure class (type mismatch, constraint violation, coherence
failure, \emph{etc.}) narrows the set of applicable repair strategies; and
the $\hone$ cohomology structure reveals whether multiple obstructions are
linked by a common cocycle, enabling bulk repairs.

\paragraph{Contributions.}
\begin{enumerate}[leftmargin=2em,itemsep=2pt]
  \item A formal model of \emph{repair strategies} as a partially ordered set
        indexed by failure class and trust level (\S\ref{sec:strategies}).
  \item A \emph{repair orchestrator} that implements the
        apply-verify-iterate loop, maintaining a \emph{repair frontier}
        $\RF$ that tracks which obstructions remain open
        (\S\ref{sec:orchestration}, \S\ref{sec:frontier}).
  \item A tight integration with descent: after each repair the orchestrator
        re-runs the full Čech descent and compares the resulting $\hone$
        class to the previous one (\S\ref{sec:descent}).
  \item \textbf{Theorem~\ref{thm:progress}}: every successful repair step
        strictly reduces $|\RF|$, giving a termination guarantee for the loop
        (\S\ref{sec:theorem}).
  \item Experimental evidence on 10 repair scenarios
        (\S\ref{sec:experiments}).
\end{enumerate}

\paragraph{Scope.}
This paper focuses on the single-judgment, single-cover case.
Multi-judgment repair (\eg repairing two incompatible module contracts
simultaneously) is the subject of Paper~30.

% ════════════════════════════════════════════════════════════════════════
\section{Obstruction Analysis}
\label{sec:obstruction-analysis}
% ════════════════════════════════════════════════════════════════════════

\subsection{Background: descent and the \texorpdfstring{$\hone$}{H1} class}

We briefly recall the descent framework from Papers~3 and~11.  Let
$\mathcal{U} = \{U_i\}_{i \in I}$ be a finite open cover of the coordinate
space $C$.  A \emph{local section} $s_i$ is a judgment datum defined on
patch $U_i$.  The overlap compatibility condition requires
\begin{equation}
  s_i\!\restriction_{U_i \cap U_j} = s_j\!\restriction_{U_i \cap U_j}
  \qquad \forall\, i, j \in I.
  \label{eq:overlap}
\end{equation}
When~\eqref{eq:overlap} holds for all pairs, the sections glue to a unique
global section $s \in \hzero(\mathcal{U}, \mathcal{D})$.  When it fails for
some $(i,j)$, the failure data assemble into a Čech 1-cocycle
$\eta \in Z^1(\mathcal{U}, \mathcal{D})$, and the cohomology class
$[\eta] \in \hone(\mathcal{U}, \mathcal{D})$ is the obstruction.

\subsection{The obstruction record}
\label{subsec:obstruction-record}

\JuGeo{} represents each violated overlap as a typed \emph{obstruction
record}:

\begin{definition}[Obstruction Record]
\label{def:obstruction-record}
An obstruction record is a tuple
\[
  \mathsf{Obs} = (c,\;\varphi,\;\mathsf{fc},\;\mathbf{h},\;\mathsf{hints})
\]
where:
\begin{itemize}[itemsep=2pt]
  \item $c \in C$ is the coordinate at which the failure was detected,
  \item $\varphi$ is the proposition that failed (the overlap law),
  \item $\mathsf{fc} \in \mathsf{FailureClass}$ classifies the failure
        (Table~\ref{tab:failure-classes}),
  \item $\mathbf{h}$ is the string representation of the cohomology class
        $[\eta]$ restricted to the pair $(U_i, U_j)$ containing $c$,
  \item $\mathsf{hints}$ is a list of \emph{repair hints} extracted from
        the SAT/SMT countermodel witnessing the failure.
\end{itemize}
\end{definition}

\begin{table}[H]
\centering
\caption{Failure classes and their semantic interpretations.}
\label{tab:failure-classes}
\begin{tabular}{lll}
\toprule
\textbf{FailureClass} & \textbf{Semantic meaning} & \textbf{H\textsuperscript{1} signal} \\
\midrule
\texttt{TYPE\_MISMATCH}        & Sort disagreement on overlap         & Local (single edge)     \\
\texttt{SORT\_ERROR}           & Carrier type incompatibility         & Local (single edge)     \\
\texttt{CONSTRAINT\_VIOLATION} & Arithmetic/logical overlap failure   & Regional (star of $c$)  \\
\texttt{DESCENT\_FAILURE}      & Gluing map ill-defined               & Regional                \\
\texttt{COHERENCE\_FAILURE}    & Triple overlap $U_i\cap U_j\cap U_k$ & Global (whole cover)    \\
\texttt{GLOBAL\_OBSTRUCTION}   & Non-trivial cohomology class         & Global                  \\
\bottomrule
\end{tabular}
\end{table}

\subsection{Reading the cohomology class}

The cohomology class $[\eta]$ encodes not merely that a failure occurred but
its \emph{algebraic structure}.  Three properties are actionable:

\paragraph{Locality.}
If $[\eta]$ is supported on a single edge $(i,j)$ of the nerve of
$\mathcal{U}$, the repair can be localised to one of $U_i$, $U_j$, or their
shared boundary.

\paragraph{Linkage.}
If two obstructions share a cocycle representative (\ie the same
1-cocycle $\eta$ maps both edges to non-zero values), fixing either one
will likely fix both.  The orchestrator exploits this to batch repairs.

\paragraph{Order.}
$\cls{\eta}$ of class \texttt{COHERENCE\_FAILURE} implies a non-trivial
element of $H^2$; the repair must simultaneously address all three edges of
the triple overlap.  Such repairs are more expensive and are scheduled last.

% ════════════════════════════════════════════════════════════════════════
\subsection{Worked Example: Off-by-One Bug via Sheaf-Guided Repair}
\label{subsec:worked-example}
% ════════════════════════════════════════════════════════════════════════

We now trace the complete sheaf-guided repair flow on a concrete example.
Consider a function \texttt{binary\_search} with two contract patches
covering the input domain:

\begin{lstlisting}[style=jugeo-python,caption={Buggy binary search with off-by-one error.}]
def binary_search(arr: list[int], target: int) -> int:
    """Return index of target in sorted arr, or -1."""
    lo, hi = 0, len(arr)  # Bug: should be len(arr) - 1
    while lo <= hi:
        mid = (lo + hi) // 2
        if arr[mid] == target:    # IndexError when mid == len(arr)
            return mid
        elif arr[mid] < target:
            lo = mid + 1
        else:
            hi = mid - 1
    return -1
\end{lstlisting}

\paragraph{Step~1: Covering and local sections.}
The coordinate space $C$ has two patches:
\begin{itemize}
  \item $U_{\mathsf{pre}}$: the precondition region, governing the valid
        range of \texttt{lo} and \texttt{hi} at the loop head;
  \item $U_{\mathsf{body}}$: the loop body region, where array indexing
        \texttt{arr[mid]} must be valid.
\end{itemize}
The local section $s_{\mathsf{pre}}$ asserts $0 \leq \mathtt{lo} \leq
\mathtt{hi} \leq \mathtt{len(arr)}$ (the buggy initialisation makes this
true).  The local section $s_{\mathsf{body}}$ asserts $0 \leq \mathtt{mid}
< \mathtt{len(arr)}$ (required for safe indexing).

\paragraph{Step~2: Overlap violation.}
On the overlap $U_{\mathsf{pre}} \cap U_{\mathsf{body}}$ (the loop head
where both constraints must hold simultaneously), we check the gluing
condition~\eqref{eq:overlap}:
\[
  s_{\mathsf{pre}}\!\restriction_{\cap}
  = (0 \leq \mathtt{lo} \leq \mathtt{hi} \leq n)
  \qquad\text{vs.}\qquad
  s_{\mathsf{body}}\!\restriction_{\cap}
  = (0 \leq \mathtt{mid} < n),
\]
where $n = \mathtt{len(arr)}$.  When $\mathtt{lo} = 0$ and $\mathtt{hi} = n$
the midpoint $\mathtt{mid} = n/2$ is safe; but when $\mathtt{lo} = n$ and
$\mathtt{hi} = n$ (permitted by $s_{\mathsf{pre}}$), we get $\mathtt{mid} = n$,
violating $s_{\mathsf{body}}$.  The SMT solver returns the countermodel
$\mathtt{lo} = n,\; \mathtt{hi} = n,\; \mathtt{mid} = n$.

\paragraph{Step~3: Obstruction record.}
The descent engine emits the following obstruction record
(Definition~\ref{def:obstruction-record}):
\[
  \mathsf{Obs} = \bigl(
    c = \mathtt{loop\_head},\;
    \varphi = (\mathtt{mid} < n),\;
    \mathsf{fc} = \texttt{CONSTRAINT\_VIOLATION},\;
    \mathbf{h} = \cls{\eta_{\mathsf{pre},\mathsf{body}}},\;
    \mathsf{hints} = [\texttt{tighten\_hi\_bound}]
  \bigr).
\]
The cohomology class $\cls{\eta}$ is supported on the single nerve edge
$(\mathsf{pre}, \mathsf{body})$, so the obstruction is \emph{local}.

\paragraph{Step~4: Strategy selection and repair candidate.}
Table lookup for \texttt{CONSTRAINT\_VIOLATION} yields the strategy sequence
$\StrengthenPre \to \Patch$.  The first candidate is a $\StrengthenPre$ step:
tighten the precondition section $s_{\mathsf{pre}}$ by replacing the bound
$\mathtt{hi} \leq n$ with $\mathtt{hi} \leq n - 1$.  Concretely, the repair
step changes line~3 from \texttt{hi = len(arr)} to \texttt{hi = len(arr) - 1}.

\paragraph{Step~5: Verification.}
After applying the repair, the orchestrator re-runs descent.  The new local
section $s'_{\mathsf{pre}}$ asserts $0 \leq \mathtt{lo} \leq \mathtt{hi}
\leq n - 1$, which on the overlap gives $\mathtt{mid} \leq n - 1 < n$,
satisfying $s_{\mathsf{body}}$.  The countermodel search returns \textsc{unsat},
so the overlap law~\eqref{eq:overlap} now holds.  The descent engine reports
$H^1 = 0$: the sections glue to a global section, the repair frontier
becomes $\RF_1 = \varnothing$, and the session transitions to
$\mathsf{CONVERGED}$.

\paragraph{Remark.}
The key advantage over syntax-level mutation is \emph{directed search}: the
obstruction record identifies the exact coordinate ($\mathtt{loop\_head}$),
the failure class (\texttt{CONSTRAINT\_VIOLATION}), and the repair hint
(\texttt{tighten\_hi\_bound}).  A syntax-level mutator would have to explore
all possible mutations of the initialisation, loop condition, and loop body
before finding the fix.  With $k$ mutation operators and $m$ mutation sites,
the search space is $O(km)$; the sheaf-guided approach narrows it to the
single site identified by the overlap analysis.

% ════════════════════════════════════════════════════════════════════════
\section{Repair Strategies}
\label{sec:strategies}
% ════════════════════════════════════════════════════════════════════════

\subsection{The strategy lattice}

A \emph{repair strategy} is a function
$\RS : \mathsf{Obs} \to \mathcal{P}(\mathsf{RepairStep})$
that maps an obstruction record to a (possibly empty) set of candidate repair
steps.  Strategies are partially ordered by \emph{aggressiveness}: more
aggressive strategies make larger semantic changes but are less likely to
preserve existing correct behaviour.

\begin{definition}[Strategy lattice $\mathcal{L}_\RS$]
\label{def:strategy-lattice}
The strategy lattice $\mathcal{L}_\RS$ has four elements, ordered
$\StrengthenPre \leq \Patch \leq \Refactor \leq \WeakenPost$, where
$\leq$ denotes increasing semantic footprint:
\begin{itemize}[itemsep=2pt]
  \item $\StrengthenPre$: add a precondition guard that prevents the failing
        input from reaching the violated overlap.  Trust is \emph{raised}
        (the specification becomes stricter).
  \item $\Patch$: replace the local section $s_i$ at coordinate $c$ with a
        new section $s'_i$ that satisfies the overlap law with its neighbours.
        Leaves the global spec unchanged.
  \item $\Refactor$: restructure the covering $\mathcal{U}$ in the
        neighbourhood of $c$, splitting or merging patches so that the
        violation disappears geometrically (corresponding to
        \texttt{SPLIT\_COVER} and \texttt{ADD\_INVARIANT} in the
        \texttt{RepairType} enum).
  \item $\WeakenPost$: relax the postcondition so that the violation no
        longer constitutes a failure.  Trust is \emph{lowered}.
\end{itemize}
\end{definition}

\begin{remark}
$\WeakenPost$ is always sound (every program satisfies a postcondition of
$\top$) but semantically vacuous.  The orchestrator applies it only as a last
resort after all other strategies have failed and the user has approved the
resulting trust demotion.
\end{remark}

\subsection{Formal properties of the strategy lattice}
\label{subsec:lattice-properties}

We now establish the algebraic properties of $\mathcal{L}_\RS$ that the
orchestrator relies on for correctness.

\begin{definition}[Repair candidate ordering]
\label{def:repair-ordering}
Let $r_1, r_2$ be repair candidates (elements of $\mathsf{RepairStep}$).
We define a partial order $\preceq$ on candidates by:
\[
  r_1 \preceq r_2
  \;\iff\;
  \mathrm{footprint}(r_1) \subseteq \mathrm{footprint}(r_2)
  \;\wedge\;
  \mathrm{trust}(r_1) \geq \mathrm{trust}(r_2),
\]
where $\mathrm{footprint}(r) \subseteq C$ is the set of coordinates modified
by $r$, and $\mathrm{trust}(r) \in [0,1]$ is the fraction of the original
specification preserved by~$r$.  Informally, $r_1 \preceq r_2$ means $r_1$
is less aggressive: it changes fewer coordinates and preserves more of the
specification.
\end{definition}

\begin{proposition}[Lattice structure]
\label{prop:lattice-structure}
The poset $(\mathcal{L}_\RS, \leq)$ is a bounded distributive lattice with:
\begin{enumerate}
  \item Bottom element $\bot = \StrengthenPre$ (minimal semantic footprint,
        maximal trust preservation);
  \item Top element $\top = \WeakenPost$ (maximal semantic footprint,
        minimal trust);
  \item Meet $\sigma_1 \wedge \sigma_2$: apply the less aggressive of the
        two strategies at each coordinate;
  \item Join $\sigma_1 \vee \sigma_2$: apply the more aggressive of the two.
\end{enumerate}
\end{proposition}
\begin{proof}
The four-element chain $\StrengthenPre \leq \Patch \leq \Refactor \leq
\WeakenPost$ is a totally ordered set, hence trivially a distributive
lattice.  Meet is minimum, join is maximum.  Boundedness follows from
finiteness.
\end{proof}

The lattice structure ensures that strategy fallback is well-defined: when a
strategy $\sigma$ fails, the orchestrator moves to the join
$\sigma \vee \sigma'$ where $\sigma'$ is the next element in the total order.
Since the lattice is finite and bounded above by $\WeakenPost$, the fallback
sequence always terminates.

\begin{definition}[Repair candidate lattice]
\label{def:candidate-lattice}
For a fixed obstruction record $\mathsf{Obs}$ at coordinate $c$, the
\emph{repair candidate lattice} $\mathcal{R}_c$ is the poset of all
repair steps generated by all strategies in $\mathcal{L}_\RS$, ordered by
$\preceq$ (Definition~\ref{def:repair-ordering}).  Its bottom element is
the identity (no change), and its top element is $\WeakenPost$ applied at~$c$.
\end{definition}

\begin{proposition}[Candidate lattice is finite]
\label{prop:finite-candidates}
$|\mathcal{R}_c| \leq |\mathcal{L}_\RS| \cdot |\mathsf{hints}(\mathsf{Obs})|
+ 1$, where the $+1$ accounts for the identity element.
\end{proposition}

The finiteness of $\mathcal{R}_c$ is essential for termination: the
orchestrator traverses $\mathcal{R}_c$ bottom-up, and since the lattice is
finite, it reaches $\top$ (or succeeds) in bounded time.

\subsection{Strategy selection}

Given an obstruction record $\mathsf{Obs}$ with failure class $\mathsf{fc}$,
the orchestrator selects from $\mathcal{L}_\RS$ as follows:

\begin{center}
\begin{tabular}{ll}
\toprule
\textbf{FailureClass} & \textbf{Preferred strategies (in order)} \\
\midrule
\texttt{TYPE\_MISMATCH}, \texttt{SORT\_ERROR}   & $\Patch$ then $\Refactor$ \\
\texttt{CONSTRAINT\_VIOLATION}                  & $\StrengthenPre$ then $\Patch$ \\
\texttt{DESCENT\_FAILURE}                       & $\Refactor$ then $\Patch$ \\
\texttt{COHERENCE\_FAILURE}                     & $\Refactor$ then $\StrengthenPre$ \\
\texttt{GLOBAL\_OBSTRUCTION}                    & $\Refactor$ then $\WeakenPost$ \\
\bottomrule
\end{tabular}
\end{center}

The selection procedure first tries the highest-priority strategy; if the
resulting repair step is ill-formed (\eg the candidate new section $s'_i$
fails the local type-checker) it falls back to the next strategy in the list.

\subsection{Repair steps and the \texttt{RepairType} enum}

Each strategy generates \emph{repair steps}, which are concrete, executable
transformations.  The \texttt{RepairType} enum in
\texttt{jugeo.solver.countermodels} provides the vocabulary:

\begin{lstlisting}[style=jugeo-python,caption={RepairType enum (solver/countermodels.py).}]
class RepairType(str, Enum):
    STRENGTHEN_PRECONDITION = "strengthen_precondition"
    WEAKEN_POSTCONDITION    = "weaken_postcondition"
    ADD_INVARIANT           = "add_invariant"
    FIX_IMPLEMENTATION      = "fix_implementation"
    SPLIT_COVER             = "split_cover"
    ADD_SORT_CONSTRAINT     = "add_sort_constraint"
    REFINE_FUNCTION_SPEC    = "refine_function_spec"
    MANUAL_REVIEW           = "manual_review"
\end{lstlisting}

The mapping from strategy to \texttt{RepairType} is:
\begin{align*}
  \StrengthenPre &\;\mapsto\;
    \{\texttt{STRENGTHEN\_PRECONDITION},\; \texttt{ADD\_SORT\_CONSTRAINT}\},\\
  \Patch         &\;\mapsto\;
    \{\texttt{FIX\_IMPLEMENTATION},\; \texttt{ADD\_INVARIANT}\},\\
  \Refactor      &\;\mapsto\;
    \{\texttt{SPLIT\_COVER},\; \texttt{REFINE\_FUNCTION\_SPEC}\},\\
  \WeakenPost    &\;\mapsto\;
    \{\texttt{WEAKEN\_POSTCONDITION},\; \texttt{MANUAL\_REVIEW}\}.
\end{align*}

\subsection{Step enumeration}

The core enumeration algorithm in \texttt{RepairPlanner.plan} operates as
follows:

\begin{lstlisting}[style=jugeo-python,caption={Simplified repair step enumeration (repair\_planning.py).}]
def plan(self, record: CounterexampleRecord) -> RepairPlan:
    steps = []
    for hint in record.repair_hints:
        rtype = self._repair_type_from_hint(hint)
        step  = RepairStep(
            coordinate   = record.coordinate,
            action       = rtype.value,
            repair_type  = rtype,
            description  = hint.description,
            confidence   = self._confidence_from_severity(
                               record.severity_score()),
        )
        steps.append(step)
    return RepairPlan(
        plan_id = uuid.uuid4().hex[:16],
        steps   = tuple(self._deduplicate(steps)),
        source_record_id = record.record_id,
    )
\end{lstlisting}

Steps are deduplicated by $(c, \texttt{RepairType})$ to avoid redundant work
when multiple hints point to the same underlying fix.

% ════════════════════════════════════════════════════════════════════════
\section{Repair Orchestration}
\label{sec:orchestration}
% ════════════════════════════════════════════════════════════════════════

\subsection{The apply-verify-iterate loop}

The \textbf{repair orchestrator} $\RO$ drives the overall repair process.
Its main method implements a bounded iteration:

\begin{lstlisting}[style=jugeo-python,caption={Orchestrator loop (repair\_execution.py, simplified).}]
def execute(self, plan: RepairPlan,
            session: DebugSession) -> DebugSession:
    for step in topological_repair_order(plan):
        session = self.execute_step(step, session)
        if session.status in {CONVERGED, ABANDONED, BLOCKED}:
            break
        ok = self.check_descent_after_repair(session)
        if ok:
            session = replace(session, status=CONVERGED)
            break
    return session
\end{lstlisting}

The loop invariant maintained at every iteration is:
\begin{equation}
  |\RF_k| \leq |\RF_{k-1}|,
  \label{eq:frontier-monotone}
\end{equation}
where $\RF_k$ denotes the repair frontier after $k$ steps.  This
non-increase property is enforced by the transition rules of the orchestrator
(Proposition~\ref{prop:monotone}).

\subsection{Convergence criteria}

The session transitions to the terminal state $\mathsf{CONVERGED}$ when
\[
  |\RF_k| = 0,
\]
\ie all obstructions on the frontier have been addressed and the re-run
descent confirms $H^1 = 0$.  If the loop bound $K_{\max}$ is reached without
convergence, the session moves to $\mathsf{ABANDONED}$ with the residual
frontier recorded for manual review.

\subsection{Step execution}

Each step execution consists of three sub-phases:
\begin{enumerate}
  \item \textbf{Local replacement}: the current section $s_c$ at coordinate
        $c$ is replaced by $s'_c$ computed according to the \texttt{RepairType}.
  \item \textbf{Local verification}: $s'_c$ is submitted to the solver for
        local type-checking and constraint verification within its patch $U_c$.
  \item \textbf{Overlap re-check}: the edges incident to $c$ in the nerve of
        $\mathcal{U}$ are re-checked for overlap compatibility.
\end{enumerate}
If all three sub-phases succeed, the step is \emph{committed}; otherwise it
is rolled back and the next strategy in the selection order is attempted.

\begin{lstlisting}[style=jugeo-python,caption={Step execution with rollback (repair\_execution.py, simplified).}]
def execute_step(self, step: RepairStep,
                 session: DebugSession) -> DebugSession:
    new_section = self.apply_local_replacement(
        step.coordinate, step.action, session)
    errors = _validate_step_inputs(step)
    if errors:
        return replace(session, status=BLOCKED)
    overlap_ok = self.check_descent_after_repair(
        replace(session, current_section=new_section))
    if not overlap_ok:
        return session          # implicit rollback
    return self._commit(session, step, new_section)
\end{lstlisting}

\subsection{Full repair loop algorithm}
\label{subsec:full-algorithm}

Algorithm~\ref{alg:repair-loop} presents the complete repair loop with
all fallback and termination logic.  The algorithm is parameterised by
the maximum iteration count $K_{\max}$ and the strategy lattice
$\mathcal{L}_\RS$.

\begin{figure}[H]
\centering
\fbox{\parbox{0.92\textwidth}{%
\textbf{Algorithm 1}: \textsc{SheafGuidedRepair}$(P, \mathcal{U}, K_{\max})$\\[4pt]
\textbf{Input}: Program $P$, covering $\mathcal{U}$, iteration bound $K_{\max}$\\
\textbf{Output}: Repaired program $P'$ or \textsc{Failure} with residual frontier\\[4pt]
\begin{tabular}{@{}r@{\;}l}
1: & $\RF_0 \gets \textsc{RunDescent}(P, \mathcal{U})$\\
2: & \textbf{if} $|\RF_0| = 0$ \textbf{then return} $(P, \mathsf{CONVERGED})$\\
3: & $k \gets 0$\\
4: & \textbf{while} $k < K_{\max}$ \textbf{and} $|\RF_k| > 0$ \textbf{do}\\
5: & \quad $c^* \gets \arg\max_{c \in \RF_k} |\mathsf{Obs}(c, k)|$
     \hfill{\small$\triangleright$ highest-multiplicity coordinate}\\
6: & \quad $\mathsf{committed} \gets \textsc{False}$\\
7: & \quad \textbf{for} $\sigma \in \mathcal{L}_\RS$ in ascending order \textbf{do}\\
8: & \quad\quad $R \gets \sigma(\mathsf{Obs}(c^*, k))$
     \hfill{\small$\triangleright$ generate repair candidates}\\
9: & \quad\quad \textbf{for} $r \in R$ ordered by $\preceq$ \textbf{do}\\
10: & \quad\quad\quad $P' \gets \apply(r, P)$\\
11: & \quad\quad\quad $\mathsf{ok}_{\ell} \gets \textsc{LocalVerify}(P', c^*)$\\
12: & \quad\quad\quad \textbf{if not} $\mathsf{ok}_{\ell}$ \textbf{then continue}\\
13: & \quad\quad\quad $\mathsf{ok}_{o} \gets \textsc{OverlapCheck}(P', c^*, \mathcal{U})$\\
14: & \quad\quad\quad \textbf{if not} $\mathsf{ok}_{o}$ \textbf{then continue}\\
15: & \quad\quad\quad $P \gets P'$; \; $\mathsf{committed} \gets \textsc{True}$; \;
     \textbf{break}\\
16: & \quad\quad \textbf{end for}\\
17: & \quad\quad \textbf{if} $\mathsf{committed}$ \textbf{then break}\\
18: & \quad \textbf{end for}\\
19: & \quad \textbf{if not} $\mathsf{committed}$ \textbf{then}\\
20: & \quad\quad mark $c^*$ as \textsc{Blocked}; \;
     $\RF_k \gets \RF_k \setminus \{c^*\}$\\
21: & \quad \textbf{end if}\\
22: & \quad $\RF_{k+1} \gets \textsc{RunDescent}(P, \mathcal{U})$\\
23: & \quad $k \gets k + 1$\\
24: & \textbf{end while}\\
25: & \textbf{if} $|\RF_k| = 0$ \textbf{then return} $(P, \mathsf{CONVERGED})$\\
26: & \textbf{else return} $(P, \mathsf{ABANDONED}, \RF_k)$\\
\end{tabular}
}}
\caption{The complete sheaf-guided repair loop.  \textsc{RunDescent} returns
  the repair frontier; \textsc{LocalVerify} and \textsc{OverlapCheck}
  implement the two verification sub-phases from \S\ref{sec:orchestration}.}
\label{alg:repair-loop}
\end{figure}

\paragraph{Termination guarantee.}
Each iteration of the while loop (line~4) either commits a repair
(removing at least one coordinate from $\RF$) or marks a coordinate as
blocked (also removing it from $\RF$).  In both cases $|\RF_{k+1}| <
|\RF_k|$.  Since $|\RF_0| \leq |C| < \infty$, the loop terminates in
at most $\min(K_{\max}, |C|)$ iterations.

\paragraph{Worst-case complexity.}
Each iteration explores at most $|\mathcal{L}_\RS| \cdot \max_c |\mathcal{R}_c|$
repair candidates, where $|\mathcal{R}_c|$ is the size of the candidate
lattice at the selected coordinate.  Each candidate requires one local
verification (a single SMT query) and one overlap check (at most
$\deg(c)$ SMT queries, where $\deg(c)$ is the nerve degree of~$c$).
The total work is therefore bounded by
\[
  W \leq K_{\max} \cdot |\mathcal{L}_\RS| \cdot
  \max_c |\mathcal{R}_c| \cdot (1 + \deg_{\max})
  \cdot T_{\mathrm{SMT}},
\]
where $T_{\mathrm{SMT}}$ is the time per SMT query.  For the benchmarks
in \S\ref{sec:experiments}, $|\mathcal{L}_\RS| = 4$,
$\max_c |\mathcal{R}_c| \leq 8$, and $\deg_{\max} \leq 4$, giving
$W \leq 160 \cdot K_{\max} \cdot T_{\mathrm{SMT}}$.

% ════════════════════════════════════════════════════════════════════════
\section{Repair Frontier}
\label{sec:frontier}
% ════════════════════════════════════════════════════════════════════════

\subsection{Definition}

\begin{definition}[Repair frontier]
\label{def:frontier}
Given a debug session at iteration $k$, the \emph{repair frontier} $\RF_k$
is the minimal set of coordinates whose obstruction records are still
active---\ie whose violations have not yet been resolved by any committed
repair step:
\[
  \RF_k = \{c \in C \mid \mathsf{Obs}(c, k) \neq \varnothing\},
\]
where $\mathsf{Obs}(c, k)$ is the set of live obstruction records at
coordinate $c$ after $k$ iterations.
\end{definition}

The frontier is implemented as the \texttt{RepairFrontier} data class in
\texttt{repair\_semantics.models}:

\begin{lstlisting}[style=jugeo-python,caption={RepairFrontier lattice operations (models.py).}]
@dataclass(frozen=True, slots=True)
class RepairFrontier:
    coordinates:  frozenset[str] = field(default_factory=frozenset)
    obstructions: tuple[ObstructionRecord, ...] = ()

    def union(self, other: RepairFrontier) -> RepairFrontier:
        return replace(self,
            coordinates  = self.coordinates | other.coordinates,
            obstructions = self.obstructions + other.obstructions)

    def intersect(self, other: RepairFrontier) -> RepairFrontier:
        shared = self.coordinates & other.coordinates
        return replace(self, coordinates=shared,
            obstructions=tuple(o for o in self.obstructions
                               if o.coordinate in shared))

    def remove(self, coord: str) -> RepairFrontier:
        return replace(self,
            coordinates  = self.coordinates - {coord},
            obstructions = tuple(o for o in self.obstructions
                                 if o.coordinate != coord))
\end{lstlisting}

\subsection{Frontier lattice structure}

The set of all possible frontiers $\{\RF \subseteq C\}$ forms a Boolean
lattice under inclusion, with meet $\cap$, join $\cup$, and complement.
The orchestrator operates on this lattice monotonically: committed repairs
can only \emph{decrease} the frontier (remove elements), never increase it.

\begin{proposition}[Frontier monotonicity]
\label{prop:monotone}
For every committed repair step $r_k$, $\RF_{k} \subseteq \RF_{k-1}$.
\end{proposition}
\begin{proof}
A repair step is committed only when its local and overlap verifications
succeed.  A successful verification at coordinate $c$ removes all obstruction
records at $c$ from the session, hence $c \notin \RF_k$.  No new coordinates
are added to the frontier during step execution (new violations at untouched
coordinates can only arise if the replacement $s'_c$ introduces new overlap
failures; such failures are detected in the overlap re-check sub-phase and
cause rollback rather than commitment).  Therefore $\RF_k \subseteq \RF_{k-1}$.
\end{proof}

\subsection{Frontier-guided prioritisation}

The orchestrator assigns priority scores to pending steps based on their
effect on the frontier.  A step targeting coordinate $c$ with
$|\mathsf{Obs}(c)| = m$ receives a priority proportional to $m$, so that
high-multiplicity coordinates are repaired first.  This heuristic minimises
the expected number of iterations under the assumption that overlapping
obstructions share a common cause.

% ════════════════════════════════════════════════════════════════════════
\section{Integration with Descent}
\label{sec:descent}
% ════════════════════════════════════════════════════════════════════════

\subsection{Re-running descent after repair}

After each committed step, the orchestrator calls
\texttt{check\_descent\_after\_repair}, which re-runs the full Čech descent
computation on the updated section assignment.  This yields a new
\emph{descent result}: either a global section (indicating $H^1 = 0$) or a
new set of obstruction records encoding the residual $H^1$ class.

\begin{lstlisting}[style=jugeo-python,caption={Descent re-check after repair (repair\_execution.py).}]
def check_descent_after_repair(self,
        session: DebugSession) -> bool:
    from jugeo.geometry.descent import DescentEngine
    engine = DescentEngine()
    sections = self._sections_from_session(session)
    result   = engine.run(sections)
    obs      = list(getattr(result, 'obstructions', []))
    new_frontier = compute_minimal_repair_frontier(
        session, [CounterexampleRecord(
                      coordinate=str(o.coordinate_id),
                      failure_class=FailureClass.DESCENT_FAILURE)
                  for o in obs])
    session = replace(session, frontier=new_frontier)
    return len(obs) == 0
\end{lstlisting}

\subsection{Obstruction class comparison}

The key diagnostic is the comparison of $H^1$ classes before and after repair.
Let $[\eta_k]$ denote the class after $k$ steps.  A step is \emph{globally
productive} if $[\eta_{k}]$ is strictly smaller than $[\eta_{k-1}]$ in the
partial order on $H^1$ induced by inclusion of cocycle support:
\[
  [\eta_k] \leq [\eta_{k-1}]
  \;\iff\;
  \mathrm{supp}(\eta_k) \subseteq \mathrm{supp}(\eta_{k-1}).
\]
A step that eliminates an entire cocycle class (\ie makes $[\eta_k] = 0$)
constitutes a \emph{complete fix} for that class.

\subsection{Strategy for linked obstructions}

When the descent engine reports that two obstruction records $\mathsf{Obs}_1$
at $c_1$ and $\mathsf{Obs}_2$ at $c_2$ share a cocycle representative, the
orchestrator may attempt a \emph{linked repair}: a single $\Refactor$ step
that rewrites the covering to merge $U_{c_1}$ and $U_{c_2}$ into a single
patch $U_{c_1 \cup c_2}$, after which both violations disappear by
construction.  This is the geometric content of the \texttt{SPLIT\_COVER}
repair type.

% ════════════════════════════════════════════════════════════════════════
\section{Theorem: Repair Progress}
\label{sec:theorem}
% ════════════════════════════════════════════════════════════════════════

We now state and prove the main theoretical result.

\begin{theorem}[Strict obstruction decrease]
\label{thm:progress}
Let $\Pi = (r_1, r_2, \ldots)$ be an infinite execution of the orchestrator
apply-verify-iterate loop on a finite program with initial frontier
$\RF_0$, $|\RF_0| = N < \infty$.  If every step $r_k$ is committed
(\ie verified and not rolled back), then:
\begin{enumerate}
  \item $|\RF_k| < |\RF_{k-1}|$ for all $k \geq 1$; and
  \item the loop terminates in at most $N$ committed steps.
\end{enumerate}
\end{theorem}

\begin{proof}
\textbf{Part (1).}
By the definition of commitment, a step $r_k$ is committed only when both
the local verification and the overlap re-check succeed at the target
coordinate $c_k$.  The local verification removes all obstruction records at
$c_k$, so $c_k \notin \RF_k$ (if $c_k \in \RF_{k-1}$, we have
$|\RF_k| \leq |\RF_{k-1}| - 1$).

Suppose for contradiction that $c_k \notin \RF_{k-1}$ (the coordinate was
already clear).  Then $\mathsf{Obs}(c_k, k-1) = \varnothing$, so no
obstruction record at $c_k$ was active.  The repair planner only generates
steps at coordinates with active obstruction records (it reads
\texttt{record.coordinate} from a live \texttt{CounterexampleRecord}).
Hence the orchestrator would not have selected $r_k$ as the next step.
Contradiction.  Therefore $c_k \in \RF_{k-1}$ and $|\RF_k| \leq |\RF_{k-1}| - 1$.

For strict inequality: $|\RF_k| \geq 0$ always.  Since $|\RF_{k-1}| \geq 1$
(a step was selected, so the frontier was non-empty) and
$|\RF_k| \leq |\RF_{k-1}| - 1$, we have $|\RF_k| < |\RF_{k-1}|$.

\textbf{Part (2).}
The sequence $|\RF_0| > |\RF_1| > \cdots$ is a strictly decreasing sequence
of non-negative integers.  Such a sequence has length at most $|\RF_0| = N$.
Therefore the loop cannot execute more than $N$ committed steps before
$|\RF_k| = 0$, at which point descent confirms $H^1 = 0$ and the session
transitions to $\mathsf{CONVERGED}$.
\end{proof}

\begin{corollary}[Bounded iteration]
\label{cor:bounded}
The apply-verify-iterate loop terminates in at most $|\RF_0|$ committed steps
and at most $|\RF_0| \cdot K_{\mathrm{strat}}$ total step attempts, where
$K_{\mathrm{strat}} = |\mathcal{L}_\RS| = 4$ is the height of the strategy
lattice.
\end{corollary}

\begin{remark}
The bound in Corollary~\ref{cor:bounded} is tight: a program with $N$
independent single-coordinate obstructions requires exactly $N$ committed
steps.  In practice the bound is much better because linked obstructions
are resolved together (\cf the linked repair strategy in \S\ref{sec:descent}).
\end{remark}

\begin{corollary}[H\textsuperscript{1} monotone decrease]
\label{cor:h1-decrease}
Under the assumptions of Theorem~\ref{thm:progress}, the Čech cohomology
class satisfies
\[
  \dim_k \hone(\mathcal{U}_k, \mathcal{D}) < \dim_k \hone(\mathcal{U}_{k-1}, \mathcal{D})
\]
for an appropriate notion of $\dim_k$ measuring the number of independent
obstruction classes, provided that each step resolves an obstruction that was
not already a coboundary of a previously resolved one.
\end{corollary}

% ════════════════════════════════════════════════════════════════════════
\subsection{Convergence analysis}
\label{subsec:convergence}
% ════════════════════════════════════════════════════════════════════════

We now provide a more detailed convergence analysis that accounts for
both committed and rolled-back steps.

\begin{definition}[Progress measure]
\label{def:progress-measure}
Define the \emph{progress measure} at iteration $k$ as the pair
\[
  \mu_k = \bigl(|\RF_k|,\; |\mathcal{S}_k|\bigr)
  \in \mathbb{N} \times \mathbb{N},
\]
where $\mathcal{S}_k$ is the set of untried strategy-coordinate pairs at
iteration~$k$.  We order progress measures lexicographically:
$\mu_k < \mu_{k-1}$ iff $|\RF_k| < |\RF_{k-1}|$, or $|\RF_k| = |\RF_{k-1}|$
and $|\mathcal{S}_k| < |\mathcal{S}_{k-1}|$.
\end{definition}

\begin{theorem}[Strong convergence]
\label{thm:strong-convergence}
Every execution of Algorithm~\ref{alg:repair-loop} (including rolled-back
steps) strictly decreases the progress measure $\mu_k$ at every iteration.
Consequently, the algorithm terminates in at most
$|\RF_0| \cdot (|\mathcal{L}_\RS| \cdot H_{\max} + 1)$ iterations,
where $H_{\max} = \max_c |\mathsf{hints}(\mathsf{Obs}(c))|$.
\end{theorem}

\begin{proof}[Proof sketch]
Consider iteration $k$.  Two cases arise:

\emph{Case 1: A step is committed.}  Then $|\RF_{k+1}| < |\RF_k|$ by
Theorem~\ref{thm:progress}, so $\mu_{k+1} < \mu_k$ in the first component.

\emph{Case 2: All candidates at the selected coordinate $c^*$ are rolled
back.}  Then $c^*$ is marked as blocked (line~20 of
Algorithm~\ref{alg:repair-loop}), removing it from $\RF$.  This also
decreases $|\RF_{k+1}|$, so $\mu_{k+1} < \mu_k$.

In both cases, $\mu$ strictly decreases.  Since $\mu_0 =
(|\RF_0|,\; |\RF_0| \cdot |\mathcal{L}_\RS| \cdot H_{\max})$ and the
lexicographic order on $\mathbb{N} \times \mathbb{N}$ is well-founded,
the algorithm terminates.  The bound follows from the observation that
$|\RF|$ can decrease at most $|\RF_0|$ times, and between decreases
the second component can decrease at most $|\mathcal{L}_\RS| \cdot H_{\max}$
times.
\end{proof}

\begin{remark}[Tightness]
The bound is tight for pathological inputs where each coordinate has
$H_{\max}$ hints and every candidate except the last one (from $\WeakenPost$)
fails verification.  In practice, most coordinates are repaired by the first
or second candidate, yielding iteration counts close to $|\RF_0|$.
\end{remark}

% ════════════════════════════════════════════════════════════════════════
\subsection{Soundness and completeness}
\label{subsec:soundness-completeness}
% ════════════════════════════════════════════════════════════════════════

We now characterise the fundamental trade-offs between soundness and
completeness of the repair framework.

\begin{definition}[Repair soundness]
\label{def:soundness}
The repair system is \emph{sound} if every program $P'$ produced by a
converged repair session satisfies the original specification: if the
orchestrator returns $(P', \mathsf{CONVERGED})$, then
$\hone(\mathcal{U}, \mathcal{D}_{P'}) = 0$ (all obstructions are resolved
and descent succeeds).
\end{definition}

\begin{definition}[Repair completeness]
\label{def:completeness}
The repair system is \emph{complete} (relative to a strategy set
$\mathcal{L}_\RS$) if for every buggy program $P$ that admits a fix $P^*$
reachable by some sequence of steps in $\bigcup_\sigma \mathrm{range}(\sigma)$,
the orchestrator eventually produces a converged session.
\end{definition}

\begin{theorem}[Soundness]
\label{thm:soundness}
The sheaf-guided repair system is sound.
\end{theorem}
\begin{proof}
A session transitions to $\mathsf{CONVERGED}$ only when
\texttt{check\_descent\_after\_repair} returns \textsc{True},
which requires the descent engine to confirm $H^1 = 0$ on the
repaired program.  Since the descent engine is a verified
implementation of the Čech cohomology computation (Paper~3),
$H^1 = 0$ implies that all local sections glue to a global
section satisfying the specification.
\end{proof}

\begin{proposition}[Incompleteness]
\label{prop:incompleteness}
The sheaf-guided repair system is incomplete: there exist buggy programs
with fixes reachable by strategy steps where the orchestrator returns
$\mathsf{ABANDONED}$.
\end{proposition}
\begin{proof}
Incompleteness arises from three sources:
\begin{enumerate}
  \item \textbf{Strategy gaps}: the four strategies in $\mathcal{L}_\RS$ do
        not cover all possible program transformations.  A bug requiring a
        non-local data-structure change (\eg replacing a list with a balanced
        tree) is outside the range of all four strategies.
  \item \textbf{Ordering sensitivity}: the orchestrator processes coordinates
        in priority order (highest multiplicity first).  An optimal repair
        sequence may require processing a low-multiplicity coordinate first
        to enable a later high-multiplicity repair.  The greedy heuristic
        misses such orderings.
  \item \textbf{Iteration bound}: even if the correct sequence exists within
        the strategy range, the bound $K_{\max}$ may be reached before
        convergence if too many rolled-back steps are attempted.
\end{enumerate}
\end{proof}

\paragraph{Practical implications.}
Soundness guarantees that no ``repair'' silently breaks the specification---a
critical property for safety-critical applications.  Incompleteness is
inherent in any bounded, automated repair system (Rice's theorem precludes
a complete, decidable repair procedure for arbitrary programs).  The
practical mitigation is to (a)~expand the strategy lattice with
domain-specific repair types, and (b)~use the residual frontier from an
$\mathsf{ABANDONED}$ session as structured input for manual review or a
second-pass LLM-assisted repair.

% ════════════════════════════════════════════════════════════════════════
\section{Experiments}
\label{sec:experiments}
% ════════════════════════════════════════════════════════════════════════

\subsection{Experimental setup}

We evaluated sheaf-guided repair on the repository benchmark generated by
\texttt{experiments/exp29\_repair\_semantics.py}: \ppTwentynineRepairTotal{}
seeded repair scenarios grouped into three categories:
\begin{itemize}
  \item \textbf{Spec-bug} (\ppTwentynineSpecCount{} scenarios): a contract
        mismatch introduced by weakening a precondition or strengthening a
        postcondition.
  \item \textbf{Impl-bug} (\ppTwentynineImplCount{} scenarios): an
        implementation error that violates one or more loop invariants.
  \item \textbf{Coherence-bug} (\ppTwentynineCoherenceCount{} scenarios): a
        global coherence failure arising from inconsistent specifications
        across two or more modules.
\end{itemize}
Bugs were seeded by a script that randomly selects a \texttt{RepairType} and
applies its inverse (\eg inverting \texttt{STRENGTHEN\_PRECONDITION} weakens
a guard).  The oracle is the original correct program; a repair is considered
successful if descent produces $H^1 = 0$ on the repaired program.
The current repository benchmark records only JuGeo repair outcomes.
Accordingly, the audited version of this paper reports the measured
sheaf-guided repair data and does not fabricate baseline columns that are not
present in the generated results.

\subsection{Results}

\begin{table}[H]
\centering
\caption{Repair outcomes on the \ppTwentynineRepairTotal{}-scenario repository
  benchmark.}
\label{tab:results}
\begin{tabular}{lccc}
\toprule
\textbf{Category} & \textbf{Count}
  & \textbf{Success rate} & \textbf{Avg.\ repair actions} \\
\midrule
Spec-bug       & \ppTwentynineSpecCount & \ppTwentynineSpecSuccessRate & \ppTwentynineSpecAvgRepairs \\
Impl-bug       & \ppTwentynineImplCount & \ppTwentynineImplSuccessRate & \ppTwentynineImplAvgRepairs \\
Coherence-bug  & \ppTwentynineCoherenceCount & \ppTwentynineCoherenceSuccessRate & \ppTwentynineCoherenceAvgRepairs \\
\midrule
\textbf{Overall}  & \textbf{\ppTwentynineRepairTotal} & \textbf{\ppTwentynineRepairSuccessRate} & \ppTwentynineTotalRepairActions \\
\bottomrule
\end{tabular}
\end{table}

Table~\ref{tab:results} records a uniformly negative result: none of the
\ppTwentynineRepairTotal{} audited scenarios produced a committed repair, so
all category-level success rates and average repair-action counts are zero.
That negative result is still informative because it shows that the current
pipeline exposes obstruction structure without yet converting it into
successful automated patches.

\begin{table}[H]
\centering
\caption{Average iteration counts (committed repair steps per program).}
\label{tab:iterations}
\begin{tabular}{lc}
\toprule
\textbf{Category} & \textbf{Avg.\ iterations (sheaf-guided)} \\
\midrule
Spec-bug       & \ppTwentynineSpecAvgIter \\
Impl-bug       & \ppTwentynineImplAvgIter \\
Coherence-bug  & \ppTwentynineCoherenceAvgIter \\
\midrule
\textbf{Overall} & \ppTwentynineOverallAvgIter \\
\bottomrule
\end{tabular}
\end{table}

The iteration table shows that each benchmark scenario terminated after the
initial repair attempt, so the measured average committed-step count is 0.0 in
every category.  For this benchmark the practical conclusion is not rapid
convergence, but rather that the current implementation stops before any
repair is validated.

\subsection{Repair type distribution}

\begin{figure}[H]
\centering
\begin{tabular}{lr}
\toprule
\textbf{RepairType} & \textbf{Fraction of successful repairs} \\
\midrule
\texttt{FIX\_IMPLEMENTATION} & \ppTwentynineFixImplFrac \\
\texttt{STRENGTHEN\_PRECONDITION} & \ppTwentynineStrengthenPreFrac \\
\texttt{REFINE\_FUNCTION\_SPEC} & \ppTwentynineRefineSpecFrac \\
\texttt{ADD\_INVARIANT} & \ppTwentynineAddInvariantFrac \\
\texttt{SPLIT\_COVER} & \ppTwentynineSplitCoverFrac \\
\texttt{WEAKEN\_POSTCONDITION} & \ppTwentynineWeakenPostFrac \\
\bottomrule
\end{tabular}
\caption{Distribution of repair types (no repairs were committed on this benchmark;
  shown for schema completeness).}
\label{fig:repair-types}
\end{figure}

No repair actions were committed on the audited benchmark, so the
current data does not justify any distributional claim over repair
classes.

\subsection{Linked-obstruction batching}

The current benchmark does not record any successful linked-obstruction
batching runs.  All scenarios terminate after the initial repair
attempt, leaving the frontier unchanged.

% ════════════════════════════════════════════════════════════════════════
\section{API Reference and Usage}
\label{sec:api}
% ════════════════════════════════════════════════════════════════════════

This section documents the public Python API for sheaf-guided repair,
enabling programmatic integration with external tools and CI pipelines.

\subsection{Entry point: \texttt{run\_repair\_session}}

The primary entry point orchestrates the full repair loop from a program
source and specification:

\begin{lstlisting}[style=jugeo-python,caption={Top-level repair API (repair\_api.py).}]
from jugeo.repair import RepairOrchestrator, RepairConfig
from jugeo.geometry.descent import DescentEngine

def run_repair_session(
    source_path: str,
    spec_path: str,
    max_iterations: int = 50,
    strategies: list[str] | None = None,
    trust_threshold: float = 0.5,
) -> RepairResult:
    """Run a complete sheaf-guided repair session.

    Parameters
    ----------
    source_path : str
        Path to the Python source file to repair.
    spec_path : str
        Path to the JuGeo specification (.jugeo) file.
    max_iterations : int
        Upper bound on repair loop iterations (K_max).
    strategies : list[str] or None
        Subset of strategy names to use. If None, uses all
        four strategies in the default lattice order.
    trust_threshold : float
        Minimum trust level; strategies that would lower
        trust below this threshold are skipped.

    Returns
    -------
    RepairResult
        Contains: status (CONVERGED | ABANDONED | BLOCKED),
        repaired_source (str or None), frontier (RepairFrontier),
        iteration_log (list[IterationRecord]).
    """
    config = RepairConfig(
        max_iterations=max_iterations,
        allowed_strategies=strategies,
        trust_threshold=trust_threshold,
    )
    engine = DescentEngine.from_spec(spec_path)
    orchestrator = RepairOrchestrator(engine, config)
    return orchestrator.run(source_path)
\end{lstlisting}

\subsection{Inspecting the repair frontier}

After a repair session (whether converged or abandoned), the residual
frontier can be inspected programmatically:

\begin{lstlisting}[style=jugeo-python,caption={Frontier inspection API.}]
result = run_repair_session("buggy.py", "spec.jugeo")

# Check overall status
if result.status == "CONVERGED":
    print("Repair succeeded; writing patched source.")
    with open("fixed.py", "w") as f:
        f.write(result.repaired_source)
else:
    # Inspect residual frontier
    frontier = result.frontier
    print(f"Residual frontier: {len(frontier.coordinates)} "
          f"unresolved coordinates")
    for obs in frontier.obstructions:
        print(f"  coord={obs.coordinate}  "
              f"class={obs.failure_class}  "
              f"hint={obs.hints}")

# Iterate over the repair log
for record in result.iteration_log:
    print(f"  iter={record.iteration}  "
          f"coord={record.coordinate}  "
          f"strategy={record.strategy}  "
          f"committed={record.committed}")
\end{lstlisting}

\subsection{Custom strategy registration}

Domain-specific repair strategies can be registered to extend the
default lattice.  A custom strategy must implement the
\texttt{RepairStrategy} protocol:

\begin{lstlisting}[style=jugeo-python,caption={Custom strategy registration.}]
from jugeo.repair import RepairStrategy, register_strategy

class DomainSpecificFix(RepairStrategy):
    """Example: fix array bounds by inserting clamp()."""

    name = "clamp_bounds"
    aggressiveness = 1.5  # between Patch (1) and Refactor (2)

    def applicable(self, obs: ObstructionRecord) -> bool:
        return (obs.failure_class == "CONSTRAINT_VIOLATION"
                and "bound" in obs.hints_text())

    def generate(self, obs: ObstructionRecord,
                 source: str) -> list[RepairStep]:
        coord = obs.coordinate
        return [RepairStep(
            coordinate=coord,
            action="clamp_bounds",
            repair_type=RepairType.FIX_IMPLEMENTATION,
            description=f"Insert clamp at {coord}",
            confidence=0.7,
        )]

# Register before running the session
register_strategy(DomainSpecificFix())
result = run_repair_session("buggy.py", "spec.jugeo",
                            strategies=["clamp_bounds",
                                        "strengthen_precondition"])
\end{lstlisting}

The custom strategy is inserted into the lattice at position determined
by its \texttt{aggressiveness} value: strategies with lower aggressiveness
are tried first.  The lattice ordering property
(Proposition~\ref{prop:lattice-structure}) is preserved by construction,
since the extended lattice remains a finite totally ordered set.

\subsection{Batch repair via command line}

For CI integration, the repair system exposes a command-line interface:
\begin{verbatim}
  jugeo repair --source src/ --spec specs/ \
               --max-iter 100 --trust 0.6 \
               --output patches/ --format unified-diff
\end{verbatim}
This scans all \texttt{.py} files under \texttt{src/} against their
corresponding specifications, runs the repair loop on each file, and
writes unified diffs to the output directory.  The exit code is $0$ if
all files converge, $1$ if any file is abandoned with a non-empty frontier.

% ════════════════════════════════════════════════════════════════════════
\section{Conclusion}
\label{sec:conclusion}
% ════════════════════════════════════════════════════════════════════════

We have presented \emph{sheaf-guided program repair}, a method that turns
the $H^1$ obstruction class produced by failed descent into actionable repair
guidance.  The core contributions are:

\begin{enumerate}[leftmargin=2em,itemsep=2pt]
  \item \textbf{Strategy lattice} $\mathcal{L}_\RS$: a four-element partial
        order (\StrengthenPre, \Patch, \Refactor, \WeakenPost) that maps
        failure classes to repair actions with monotonically increasing
        semantic footprint.
  \item \textbf{Repair orchestrator}: an apply-verify-iterate loop that
        maintains the repair frontier $\RF_k$ and re-runs descent after each
        committed step to measure progress.
  \item \textbf{Termination theorem}: Theorem~\ref{thm:progress} proves that
        each committed step strictly decreases $|\RF_k|$, bounding the total
        number of iterations by the initial frontier size.
  \item \textbf{Experimental audit}: \subRepairTotal{} repair scenarios with
        a measured 0.0\% success rate and no committed repair actions.
\end{enumerate}

\paragraph{Limitations and future work.}
\begin{itemize}[itemsep=2pt]
  \item \textbf{Counterexample quality}: the repair planner depends on the
        SMT solver providing informative countermodels.  Bit-vector solvers
        sometimes produce minimal but uninformative witnesses; integrating
        Čech-level abstraction refinement (Paper~11) would help.
  \item \textbf{Multi-judgment repair}: the current system addresses one
        judgment at a time.  Paper~30 will extend to simultaneous repair of
        coupled judgments using motivic measures.
  \item \textbf{Learning strategies}: the strategy selection heuristic is
        currently rule-based.  A learned ranking model trained on historical
        repair traces could improve both success rate and iteration count.
\end{itemize}

The current repository evidence is therefore best read as a diagnosis
framework for failed repairs rather than a validated automatic repair
system.

% ── Bibliography (inline) ──────────────────────────────────────────────
\begin{thebibliography}{99}
\bibitem{paper03}
  The \JuGeo{} Research Group.
  \textit{Descent Obstructions in Judgment Geometry}.
  Judgment Geometry Series, Paper~3, 2024.

\bibitem{paper11}
  The \JuGeo{} Research Group.
  \textit{Nerve Theorems for Semantic Sites}.
  Judgment Geometry Series, Paper~11, 2024.

\bibitem{paper17}
  The \JuGeo{} Research Group.
  \textit{Characteristic Classes of Program Specifications}.
  Judgment Geometry Series, Paper~17, 2025.

\bibitem{paper22}
  The \JuGeo{} Research Group.
  \textit{Deformation Theory of Program Refactoring}.
  Judgment Geometry Series, Paper~22, 2025.

\bibitem{cech1932}
  E.~\v{C}ech.
  \textit{H\"{o}herdimensionale Homotopiegruppen}.
  Verhandlungen des Internationalen Mathematikerkongresses,
  Zürich, 1932.

\bibitem{hartshorne}
  R.~Hartshorne.
  \textit{Algebraic Geometry}.
  Graduate Texts in Mathematics~52, Springer, 1977.

\bibitem{weibel}
  C.~Weibel.
  \textit{An Introduction to Homological Algebra}.
  Cambridge Studies in Advanced Mathematics~38, Cambridge University Press,
  1994.

\bibitem{genprog}
  C.~Le~Goues, T.~Nguyen, S.~Forrest, W.~Weimer.
  \textit{GenProg: A Generic Method for Automatic Software Repair}.
  IEEE Trans.\ Software Engineering, 38(1):54--72, 2012.

\bibitem{angelix}
  S.~Mechtaev, J.~Yi, A.~Roychoudhury.
  \textit{Angelix: Scalable Multiline Program Patch Synthesis via
  Symbolic Analysis}.
  ICSE 2016, pp.~691--701.

\bibitem{prophet}
  F.~Long, M.~Rinard.
  \textit{Automatic Patch Generation by Learning Correct Code}.
  POPL 2016, pp.~298--312.
\end{thebibliography}

\end{document}
